% 机械波（高中）
% license Usr
% type Tutor

\subsection{机械波}
机械波是机械振动在空间中传播的现象。称振动的源头为波源，称承载振动的物质为介质\footnote{真空实际并非空无一物，而是能量最低的一个状态。实际上，从量子物理的观点来看，无论真空能是否为0（我们可以灵活地调整能量零点），真空总是会在极短时间内产生和湮灭粒子-反粒子对。这种涨落会被时间“抹平”，所以在相当长一段时间内观察，能量平均值还是基态。但极短时间的这种量子涨落亦能传播振动，且已经为实验所证实——\href{https://www.rmzxw.com.cn/c/2023-08-25/3399291.shtml}{真空传声}}。从微观上看，只要有波源在作振动，且质点之间有弹性恢复力，那么必然引起传播方向上，下一个相邻质点作振动。每一个质点，都可以看作接下来的，机械振动的波源，如果没有其他能量损耗，机械波的传播是没有止境的。其传播路径上的每个质点都可视为弹簧振子，在受到上一个质点的扰动后，才开始在平衡位置附近作振动。

从上述过程来看，机械波对振动能量的传播不言而喻。实际上，机械波还是振动状态（相位）的传播。假设振动在传播过程中没有损耗，如果A点为较近点，B点是较远点，且振动从A传播到B需要时间$t$，那么B在某时刻$t_0$的相位总是B在$t_0-t$的相位。所以不同质点的振动图像都有共同属性如周期，振幅，有所不同的只是相位。
\subsubsection{简谐波的描述与图像}
如果我们把机械波简化成一连串质点的运动，且设波源作简谐振动，介质是理想介质（即波速稳定），那么在任意时刻，给这一串质点“拍照”，它们的位置构成简谐函数的波形。
比如设波源的振动形式为$y=\sin(\omega t)$，波速为$v$，那么距波源$x$的质点位移为
\begin{equation}
y=\sin [\omega (t-\frac{x}{v})]~.
\end{equation}
显然如果在某时刻拍照，固定$t$后的质点位移$y$是关于$x$的简谐函数，且周期与振动周期相同，常用$T$表示。%%插图，振动和波动

我们称波在一个周期传播的距离为波长，常用$\lambda $表示，满足$\lambda=vT$。波源的振动决定周期，或者说频率，所以无论机械波经过何种介质，其周期都是不变的。然而，介质的物理性质决定了波速，所以改变介质就会改变波速和波长。
对应山峰和山谷，我们称波的位移最大值为波峰，最小值为波谷。对于简谐波，波峰和波谷的相位差总是$\pi$。

\subsubsection{波的叠加}
既然机械波实际上代表了振动在质点之间的传播，那么不同机械波在同一点处的位移之和自然是该处质点的实际位移。如果不同简谐波都满足以下条件：
\begin{enumerate}
\item 频率相同
\item 相位差恒定
\item 有沿着同一方向的振动分量
\end{enumerate}
此时，我们就称波的叠加为\textbf{波的干涉}。其特点是，两列波在任意一处的相位差恒定，任意一处的振子都在做简谐振动且振幅稳定。下面我们来证明这一点。

假设空间中有两个波源A和B，调节零时刻使得振动方程分别为
\begin{equation}
\begin{aligned}
y_A&=A\sin(\omega t)\\
y_B&=B\sin (\omega t+\phi)~,
\end{aligned}
\end{equation}
设从A,B到任意P点的时间分别为$t_1,t_2$，则在P点的合振动为

\begin{equation}
y=y_A+y_B=A\sin(\omega t-\omega t_1)+B\sin(\omega t-\omega t_2+\phi)~.
\end{equation}

显然，相位差是恒定的。

将两个波的振动比拟为匀速圆周运动后，我们可以求出合振动的形式为，

显然其振幅是稳定不变，与相位差有关的。